\begin{proposition}[Closed period groups at a rank-three fibre]\hypertarget{tag-032L}{}\label{push:src:D.4.3}
\StatementRecord{push:src:D.4.3}{proved}{032L}{\SeriesRef{fam:src:C.3.1}, \SeriesRef{fam:src:C.3.2}}

At a zero of \(D_c\), let \(A:\mathbb R^4\to\mathbb C^2\) be
the real period map and \(\ker A=\mathbb R k\), with the kernel
vector stated in \cref{fam:src:C.3.1}, \cref{fam:src:C.3.2}. Define
\[
 L_k=\{m\in\mathbb Z^4:m\cdot k=0\},\qquad r=\operatorname{rank}L_k,
 \qquad V_k=\bigcap_{m\in L_k}\ker m.
\]
Then
\[
 \overline{A\mathbb Z^4}=A(\mathbb Z^4+V_k),\qquad
 \mathbb C^2/\overline{A\mathbb Z^4}\simeq
                         \mathbb R\times(S^1)^r
\]
as real topological groups. The original quotient is a complex Lie
group precisely when \(k\) is rational, in which case \(r=3\).
Otherwise it is non-Hausdorff. Every Chabauty limit of nearby full
lattices contains \(\overline{A\mathbb Z^4}\); the quotient topology
at the fixed zero and a limit of full lattices need not coincide.
\end{proposition}


\begin{proof}
\begin{enumerate}[label=\arabic*.,leftmargin=*]
\item The last three period columns are real-linearly independent, so the
rank is exactly three. This confirms the linear algebra of \cref{fam:src:C.3.1}, \cref{fam:src:C.3.2}
without a positivity assumption. The annihilator characters of
\(A\mathbb Z^4\) are the real covectors \(\lambda\) with
\(A^t\lambda\in\mathbb Z^4\). Their integral labels belong exactly
to \(L_k\). The smallest rational subspace containing \(k\) is
\(V_k\); the closure of the image of \(\mathbb Z^4\) in
\(\mathbb R^4/\mathbb R k\) is its image after adding \(V_k\).
For an elementary verification, the two groups have the same
annihilator characters, and characters separate closed subgroups of
a real vector space by putting such a subgroup into the form
\(\mathbb R^d\times\mathbb Z^e\). This proves the closure formula.
Inside the three-dimensional image hyperplane its connected part
has dimension \(3-r\), and its discrete quotient has rank \(r\).
The one transverse real direction gives the displayed group.

\item The image is discrete exactly when the kernel line has a primitive
integral generator: extend such a generator to an integral basis for
the converse; for necessity, a discrete subgroup of a real
three-space has rank at most three, forcing a nonzero integral
kernel. A nondiscrete finitely generated subgroup cannot be closed,
since a nondiscrete closed subgroup contains an uncountable vector
group. Finally each fixed integral period converges to its limiting
period; a Chabauty limit is closed and therefore contains their
closure. This does not identify it with that closure.
\end{enumerate}
\end{proof}